\documentclass[../main.tex]{subfiles}

\begin{document}

\chapter{一般的 $L^{p}$ 空间: 完备性、对偶性和弱收敛性}\label{chapter:ch19-lp-spaces}

对于测度空间 $(X,\mathcal{M},\mu)$ 与 $1\leqslant p\leqslant\infty$, 如同在第一部分对实直线上的 Lebesgue 测度所做的那样, 我们定义线性空间 $L^{p}(X,\mu)$. 用非常类似于实直线上的 Lebesgue 测度情形的方法, 我们证明了 H\"older 与 Minkowski 不等式成立且 $L^{p}(X,\mu)$ 是 Banach 空间. 19.1 节介绍以上内容以及相关的题材. 本章的其余部分用于建立那些证明思想超出第一部分给出的思想之外的结果. 在 19.2 节, 我们用 Radon-Nikodym 定理证明关于 $L^{p}(X,\mu)$ 的对偶空间的 Riesz 表示定理, 其中 $1\leqslant p<\infty$, 而 $\mu$ 是 $\sigma$ 有限测度. 在 19.3 节, 我们证明关于 $L^{\infty}(X,\mu)$ 的对偶的 Kantorovich 表示定理. 在 19.4 节, 我们证明对于 $1<p<\infty$, Banach 空间 $L^{p}(X,\mu)$ 是自反的, 因此有自反空间具有的弱序列紧性质. 19.5 节研究非自反的 Banach 空间 $L^{1}(X,\mu)$ 的弱序列紧性. 我们用 Vitali-Hahn-Saks 定理证明 Dunford-Pettis 定理, 它告诉我们, 若 $\mu(X)<\infty$, 则 $L^{1}(X,\mu)$ 中的每个一致可积的有界序列有弱收敛的子序列.

\subfile{Chapter-19/section-01.tex}

\subfile{Chapter-19/section-02.tex}

\subfile{Chapter-19/section-03.tex}

\subfile{Chapter-19/section-04.tex}

\subfile{Chapter-19/section-05.tex}

\end{document}